% Class-compatibility preamble retained for reader excerpts of
%   Hiraku Nakajima, "Quiver Varieties and Branching", SIGMA 5 (2009) 003,
%   DOI 10.3842/SIGMA.2009.003, CC BY 4.0.
%
% The source pins the publisher class sigma, which is not part of the delivery baseline, so this
% file supplies the amsart baseline and then copies the author's own declaration block mechanically
% and unchanged: the aenume environment, the Theorem/Corollary/Lemma/Proposition and
% definition-style newtheorem declarations, the equation numbering, the author's reference macros
% (\lemref and the like), and the author's mathematical macros from the end of the preamble
% (\Gr, \Rcal, \dbar, \linf, \Irr, \Res, \Lareg, \MR, \aff and the rest).
%
% Two things are deliberately left out.  First, the publisher class internals and the single macro
% \rf that they serve, together with \graysquare and the colour \halfgray that go with it:
% nothing in the delivered excerpts uses them, and \rf depends on the internal dimension
% \y@boxdim of the sigma class.  Second, the author's package mathdots, which is not installed in
% the delivery TeX installation and whose commands no delivered excerpt uses.  Third, the author's
% no-op silencer declaration for the MR review numbers, because the delivery baseline package set
% already provides that command name and no delivered excerpt calls it.
%
% No mathematical macro, statement, symbol or assumption is introduced, changed or removed by this
% file.
\documentclass[10pt,a4paper]{amsart}
\usepackage{amssymb}
\usepackage{amsbsy,amscd,mathtools}
\usepackage{mathrsfs}
\usepackage{nicefrac}
\usepackage[all]{xy}
\usepackage{color}
\usepackage{hyperref}
\newenvironment{aenume}{%
  \begin{enumerate}%
  \renewcommand{\theenumi}{\alph{enumi}}%
  \renewcommand{\labelenumi}{\textup{(\theenumi)}}%
  }{\end{enumerate}}
%

\usepackage{verbatim}
\usepackage{version}
\usepackage{color}
\newenvironment{NB}{
\color{red}{\bf NB}. \footnotesize
}{}
% For a final version, NB disappears.
\excludeversion{NB}
% Older version
%\renewenvironment{NB}{%
%  \comment
%  }{\endcomment}
% The following is used inside NB
\newenvironment{NB2}{
\color{blue}{\bf NB}. \footnotesize
}{}
%       Theorem environments
%% \theoremstyle{plain} %% This is the default
\newtheorem{Theorem}[equation]{Theorem}
\newtheorem{Corollary}[equation]{Corollary}
\newtheorem{Lemma}[equation]{Lemma}
\newtheorem{Proposition}[equation]{Proposition}
%\newtheorem{Axiom}{Axiom}

\theoremstyle{definition}
\newtheorem{Definition}[equation]{Definition}
\newtheorem{Example}[equation]{Example}
\newtheorem{CounterExample}[equation]{Counter-Example}
\newtheorem{Notation}[equation]{Notation}
\newtheorem{Convention}[equation]{Convention}
\newtheorem{Conjecture}[equation]{Conjecture}

%\theoremstyle{remark}
\newtheorem{Remark}[equation]{Remark}
\newtheorem{Remarks}[equation]{Remarks}
%\newtheorem*{Claim}{Claim}
%\renewcommand{\theClaim}{} % to make the claim environment unnumbered
\newtheorem{Question}[equation]{Question}
%\renewcommand{\theQuestion}{}  % to make the notation environment
                               % unnumbered
\newtheorem{Exercise}[equation]{Exercise}
%\renewcommand{\theExercise}{}

\numberwithin{equation}{section}
%\numberwithin{Theorem}{section}
%\numberwithin{figure}{section}

\newcommand{\thmref}[1]{Theorem~\ref{#1}}
\newcommand{\secref}[1]{Section~\ref{#1}}
\newcommand{\lemref}[1]{Lemma~\ref{#1}}
\newcommand{\propref}[1]{Proposition~\ref{#1}}
\newcommand{\corref}[1]{Corollary~\ref{#1}}
\newcommand{\subsecref}[1]{Section~\ref{#1}}
\newcommand{\chapref}[1]{Chapter~\ref{#1}}
\newcommand{\defref}[1]{Definition~\ref{#1}}
\newcommand{\remref}[1]{Remark~\ref{#1}}
\newcommand{\remsref}[1]{Remarks~\ref{#1}}
%
%       Math definitions
%\newcommand{\lsp}[2]{{}^{#1}{#2}}
\newcommand{\lsp}[2]{{\mskip-.3mu}{}^{#1}\mskip-1mu{#2}}
\newcommand{\defeq}{\overset{\operatorname{\scriptstyle def}}{=}}
\newcommand{\interior}{\rule{6pt}{.4pt}\rule{.4pt}{6pt}} % interior
                                                         % product
\newcommand{\C}{{\mathbb C}}
\newcommand{\Z}{{\mathbb Z}}
\newcommand{\Q}{{\mathbb Q}}
\newcommand{\R}{{\mathbb R}}
% Projective spaces
\newcommand{\proj}{{\mathbb P}}
%\newcommand{\CP}{\operatorname{\C P}}
\newcommand{\CP}{\proj}
\newcommand{\RP}{\operatorname{\R P}}
% Lie Groups
\newcommand{\SL}{\operatorname{\rm SL}}
\newcommand{\SU}{\operatorname{\rm SU}}
\newcommand{\GL}{\operatorname{GL}}
\newcommand{\PGL}{\operatorname{PGL}}
\newcommand{\U}{\operatorname{\rm U}}
\newcommand{\SO}{\operatorname{\rm SO}}
\newcommand{\grpSp}{\operatorname{\rm Sp}}
% Lie algebras
\newcommand{\algsl}{\operatorname{\mathfrak{sl}}} % because \sl="slant"
\newcommand{\su}{\operatorname{\mathfrak{su}}}
\newcommand{\gl}{\operatorname{\mathfrak{gl}}}
\newcommand{\algu}{\operatorname{\mathfrak{u}}}
\newcommand{\so}{\operatorname{\mathfrak{so}}}
\newcommand{\algsp}{\operatorname{\mathfrak{sp}}}
\newcommand{\g}{{\mathfrak g}}
\newcommand{\h}{{\mathfrak h}}
% Useful symbols
\newcommand{\Spec}{\operatorname{Spec}\nolimits}
\newcommand{\Proj}{\operatorname{Proj}\nolimits}
\newcommand{\End}{\operatorname{End}}
\newcommand{\Hom}{\operatorname{Hom}}
\newcommand{\Ext}{\operatorname{Ext}}
\newcommand{\Ker}{\operatorname{Ker}}
\newcommand{\Coker}{\operatorname{Coker}}
\newcommand{\Ima}{\operatorname{Im}}
\newcommand{\coker}{\operatorname{coker}}
\newcommand{\Endskew}{\operatorname{Endskew}}
\newcommand{\im}{\mathop{\text{\rm im}}\nolimits}
%\newcommand{\dim}{\mathop{\text{\rm dim}}\nolimits}
\newcommand{\rank}{\operatorname{rank}}
\newcommand{\ind}{\mathop{\text{\rm ind}}\nolimits}
\newcommand{\tr}{\operatorname{tr}}
\newcommand{\Tr}{\operatorname{Tr}}
\newcommand{\DB}{\overline{\partial}}
\newcommand{\pd}[2]{\frac{\partial#1}{\partial#2}}
\newcommand{\deriv}[1]{\frac{d}{d#1}}
\newcommand{\id}{\operatorname{id}}
\newcommand{\ve}{\varepsilon}
%
% notations used in this paper
%
% \newcommand{\vin}[1]{\operatorname{in}(#1)} % incoming vertex
% \newcommand{\vout}[1]{\operatorname{out}(#1)} % outgoing vertex
\newcommand{\vin}[1]{\operatorname{i}(#1)} % incoming vertex
\newcommand{\vout}[1]{\operatorname{o}(#1)} % outgoing vertex
\newcommand{\bM}{{\mathbf M}} % the vector space
\newcommand{\bN}{{\mathbf N}} % the vector space
\newcommand{\M}{{\mathfrak M}} % the moduli space
\newcommand{\Mreg}{\M^{\operatorname{s}}}
\newcommand{\La}{{\mathfrak L}} % the Lagrangian variety
\newcommand{\Fi}{{\mathfrak F}} % the fixed point variety
\newcommand{\dslash}{/\!\!/} % for algebro-geometric quotient (double slash)
\newcommand{\bv}{{\mathbf v}} % the vector v
\newcommand{\bw}{{\mathbf w}} % the vector w
\newcommand{\bC}{{\mathbf C}} % the Cartan matrix
\newcommand{\bA}{{\mathbf A}} % the adjacency matrix
\newcommand{\bI}{{\mathbf I}} % the identity matrix
\newcommand{\bn}{{\mathbf n}} % the kernel of the extended Cartan matrix
\newcommand{\bu}{{\mathbf u}}
\newcommand{\be}{{\mathbf e}}
\newcommand{\bq}{{\mathbf q}}
\newcommand{\bs}{\mathbf s}
\newcommand{\bc}{\mathbf c}
\newcommand{\codim}{\operatorname{codim}} % codimension
\newcommand{\gr}{\operatorname{gr}} % grade
\newcommand{\topdeg}{{\operatorname{top}}} % top degree
\newcommand{\Uq}{{\mathbf U}_q(\mathfrak g)} % the QUE algebra
\newcommand{\Un}{{\mathbf U}_q}
\newcommand{\Um}{\Tilde{\mathbf U}_q(\mathfrak g)} % the modified
                                                 % enveloping algebra
\newcommand{\Ua}{{\mathbf U}_q(\widehat{\mathfrak g})} % the quantum affine
                                                 % algebra
\newcommand{\Uae}{{\mathbf U}_{\varepsilon}(\widehat{\mathfrak g})} %
                                                 % the quantum affine
                                                 % algebra
\newcommand{\Ul}{{\mathbf U}_q({\mathbf L}{\mathfrak g})} % the quantum
                                                      % loop algebra
\newcommand{\Uma}{\Tilde{\mathbf U}^\bw_q(\widehat{\mathfrak g})}
%                                                 the quantum affine algebra
\newcommand{\Ui}{{\mathbf U}^{\Z}_q(\mathfrak g)}
\newcommand{\Ue}{{\mathbf U}_{\varepsilon}(\mathfrak g)}
\newcommand{\Ueone}{{\mathbf U}_{\varepsilon = 1}(\mathfrak g)}
\newcommand{\Uli}{{\mathbf U}^{\Z}_q({\mathbf L}\mathfrak g)}
\newcommand{\Ule}{{\mathbf U}_{\varepsilon}({\mathbf L}\mathfrak g)}
\newcommand{\bU}{\mathbf U} % the enveloping algebra
\newcommand{\Uqm}{{\widetilde{\mathbf U}_q}} % its modified form
\newcommand{\Pa}{{\mathfrak P}} % Hecke correspondence
\newcommand{\HomE}{\operatorname{E}}
\newcommand{\HomL}{\operatorname{L}}
\newcommand{\Zreg}{Z^{\operatorname{reg}\nolimits}}
\newcommand{\hiw}{{L}} % the highest weight module
%      \interval is used to provide better spacing after a [ that
%      is used as a closing delimiter.
\newcommand{\interval}[1]{\mathinner{#1}}
\newcommand{\shfO}{\mathcal O}
\newcommand{\Mp}{\widetilde\M}
\newcommand{\Mpp}{\widehat\M}
\newcommand{\Wedge}{{\textstyle \bigwedge}}
%
\newcommand{\Zw}{Z(\bw)}
\newcommand{\Mw}{\M(\bw)}
\newcommand{\Law}{\La(\bw)}
%\newcommand{\Zw}{Z(\bullet,\bullet;\bw)}
%\newcommand{\Mw}{\M(\bullet,\bw)}
%\newcommand{\Law}{\La(\bullet,\bw)}
\newcommand{\ngpiso}{\overset{.}{=}}
\newcommand{\idl}{\mathfrak I}
\newcommand{\Ktop}{K_{\operatorname{top}}}
\newcommand{\Kztop}{K_{0,\operatorname{top}}}
\newcommand{\Kotop}{K_{1,\operatorname{top}}}

\newcommand{\qH}{\mathbb H}
\newcommand{\D}{\mathcal D}
\newcommand{\A}{\mathcal A}
\newcommand{\incarrow}{\xrightarrow{\mathrm{inclusion}}}
\newcommand{\proarrow}{\xrightarrow{\mathrm{projection}}}
%\newcommand{\oV}[1]{\overset{\circ}{V_{#1}}}
\newcommand{\oV}[1]{\breve{V}_{#1}}
\newcommand{\bV}[2]{\bar{V}^{#1}_{#2}}
\newcommand{\bB}{\bar{B}}
\newcommand{\bi}{\bar{a}}
\newcommand{\bj}{\bar{b}}
\newcommand{\bphi}{\bar{\phi}}
\newcommand{\Gr}{\operatorname{Gr}}
\newcommand{\Pp}{\mathbf P}
\newcommand{\Rcal}{\mathcal R}
\newcommand{\dbar}{\overline\partial}
\newcommand{\linf}{\ell_\infty}
\newcommand{\Irr}{\operatorname{Irr}}
\newcommand{\Res}{\operatorname{Res}}
\newcommand{\Lareg}{\mathfrak L^{\operatorname{s}}}
%\renewcommand{\MR}[1]{}
\newcommand{\aff}{\mathrm{af\/f}}
\newcommand{\affLevi}{(\g_{I_0^0})_{\aff}}


\newcommand{\hf}{\hfill}

\newcommand{\graysquare}{\color{halfgray}\blacksquare}

%\setcounter{tocdepth}{1}

%\usepackage[bookmarks=false]{hyperref}
%\usepackage[dvips,bookmarks=false]{hyperref}
%\usepackage[initials]{amsrefs}

